\documentclass[12pt]{article}

\include{../style}
\usepackage{tikz-cd}
\begin{document}

\title{Cosc 30: Discrete Mathematics}

\author{Prishita Dharampal}
\date{}
\maketitle
\vspace{-0.3in}


\begin{problem}{1}
    \begin{enumerate}
        \item[(a)]  Prove that any graph $G$ with $n$ vertices and $n$ edges must have a vertex $v$ with degree at most $2$. 
        \item[(b)] Draw a simple graph that has $10$ edges, $10$ vertices, with no vertex with degree exactly $2$. What is the largest number of degree ($\leq 1$) vertices you can have?
    \end{enumerate}
\end{problem}

\begin{solution}
    \bbni
    \begin{enumerate}
        \item[(a)] Let $V$ be the set of all vertices, note that $|V| = n$. Assume for the sake of contradiction that there exist no nodes with degree $\leq 2$. Then the minimum degree any node can have is $3$. That is 
        \[\sum_{v \in V} \deg(v) \geq 3n\]
        
        But by the Handshake Lemma, we know, 
        \[\sum_{v \in V}\deg(v) = 2n\]
        That is, 
        \[\sum_{v \in V}\deg(v) = 2n \geq 3n \]
        A contradiction! 
        That is, any graph $G$ with $n$ vertices and $n$ edges must have at least a vertex $v$ with degree at most $2$.     
        \newpage
        \item[(b)] $7$ is the largest number of degree ($\leq 1$) vertices we can have. 
        
        \begin{tikzcd}[row sep=3em, column sep=3em]
            {} & {} & \bullet \arrow[drr, -] & \bullet \arrow[dr, -] & \bullet \arrow[d, -] & \bullet \arrow[dl, -]  & \bullet \arrow[dll, -] \\
            {} & {} & {} & {} & \bullet \arrow[dl, -] \arrow[dr, -] & {} \\
            {} & {} & {} & \bullet \arrow[rr, -] & {} & \bullet & {} \\ 
            {} &  {} & \bullet \arrow[ur, -] & {} & {} & {} & \bullet \arrow[ul, -] & {}
        \end{tikzcd}
        
        
        Note there cannot exist $8$ or more such vertices because by the handshake lemma we know, that the sum of the degrees of the verticies must add up to $20$.
        Suppose $k=8$ vertices have degree $\leq 1$. They contribute at most $8$ to the degree sum, so the remaining $2$ vertices must contribute at least $20  - 8=12$, i.e. their degrees sum to at least $12$.
        Now, since these two nodes can only have $1$ edge between them (contributing $2$ to the degree sum), adding more edges would mean connecting the nodes to one of the other $8$ nodes and hence bumping up that nodes' degree by $1$. 
        
        That is, $7$ is the maximum number of degree ($\leq 1$) vertices we can have. 
    \end{enumerate}
\end{solution}

\newpage 

\begin{problem}{2}
    \begin{enumerate}
        \item [(a)] Prove that any graph containing odd-length closed walks must also contain odd cycles.
        \item[(b)] Draw a graph that contains even-length closed walks but has no even cycles.
    \end{enumerate}
\end{problem}


\begin{solution}
    \bbni
    \begin{enumerate}
        \item[(a)] Let $W$ be an odd-length closed walk in a graph $G$, and let $W'$ be the shortest odd-length walk in $G$. Then we claim that $W'$ is an odd-cycle. 
        
        We will show this by contradiction. Assume that $W'$ is not a cycle, then there is some vertex $v$ that $W'$ visits twice. Then we can break $W'$ up into two closed walks $W_1, W_2$ at $v$, such that, 
        \[|W_1| + |W_2| = |W'|\]
        Since $|W'|$ is odd, one of $|W_1|$ or $|W_2|$ is odd, and the other even. But then the odd-length closed walk is strictly shorter than $W'$, contradicting minimality!

        That is, $W'$ must be a cycle. 
        \item[(b)]   Note that the graph below has an even-length closed walk, namely, 
        \[A \rightarrow B \rightarrow C \rightarrow A \rightarrow B \rightarrow C \rightarrow A\]
        but only one cycle (which has odd-length), 
        \[A \rightarrow B \rightarrow C \rightarrow A\]
        \begin{center}
            \begin{tikzcd}[row sep=3em, column sep=3em]
            {} & A \arrow[dl, -] \arrow[dr, -] & {} \\
            B \arrow[rr, -] & {} & C \\ 
        \end{tikzcd}
        \end{center}
    \end{enumerate}
\end{solution}
\end{document}